\chapter{Persuasion, Influence \& Negotiation}\label{ch:persuasion-influence}

% Scope: evidence-based influence — field-experiment-validated principles, dual-process
%   framing, principled negotiation (BATNA/reservation/ZOPA), organizational power,
%   the trust equation. NEW (Part VII).
% Value hook: influence converts a true value proposition into closed customers (CAC down),
%   funded rounds, hires, and surplus captured (margin up) — all toward ROIC > WACC.
%   Manipulation inverts it (short-term win, trust/LTV collapse). Closes the negotiation gap.
% Draft per house style: flowing prose + example/admonition; running-SaaS worked examples.

A product that creates genuine value still fails to capture it without the ability to
persuade. The value-creation spine in this book runs from economic insight to operational
execution, but every node on that spine --- closing a customer, funding a round, hiring an
engineer, settling a supplier price --- requires a human decision to move in the right
direction. Influence is the mechanism that converts a true value proposition into a closed
sale (reducing CAC, as in \cref{eq:cac}), a signed term sheet, and captured surplus
(widening margin and pushing ROIC above WACC, per \cref{eq:roic}). When influence
techniques are instead used to manufacture a sale that the underlying product cannot
sustain, the short-term conversion is reversed in the medium term through churn that
collapses customer lifetime value (see \cref{eq:clv}). This chapter builds the
evidence-based toolkit that separates legitimate persuasion from manipulation, and
principled negotiation from zero-sum positional bargaining.


\section{The Architecture of Influence}\label{sec:influence-principles}
% complexity: 3

Why do some messages move people while others dissolve into noise? The question has a
structural answer grounded in field experiments, not anecdote. The popular treatment
of this domain, \textcite{cialdini2021influence}, catalogues seven principles of compliance;
the scientific foundation beneath those principles is a body of randomised field
experiments that measure actual compliance behaviour in real environments, with
pre-registered control conditions and quantified effect sizes.

The most direct experimental evidence concerns descriptive norms --- the signal that
``people like you, in your situation, behave this way.'' \textcite{goldstein2008field}
conducted a field experiment inside a hotel, randomising towel-reuse messages across
occupied rooms without guests' awareness. The standard environmental appeal (``Help save
the environment'') produced a reuse rate of \qty{35.1}{\percent}. A general descriptive
norm (``The majority of guests reuse their towels'') raised compliance to
\qty{44.1}{\percent} --- a \qty{9}{\percent}point absolute lift, representing a roughly
\qty{25.6}{\percent} relative increase over the control.
% goldstein2008field: descriptive-norm lift vs control, hotel-towel field experiment
A second condition tightened the norm to the individual's immediate context: ``\qty{75}{\percent}
of guests who stayed in this room have reused their towels.'' This provincial norm lifted
compliance further to \qty{49.3}{\percent}.
% goldstein2008field: provincial-norm rate, second experiment
The mechanism is the same one that governs conversion-rate optimisation: people use
the observed behaviour of a reference group to resolve uncertainty about the correct
action, and the more closely the reference group mirrors their own situation, the
stronger the pull.

Reciprocity has been measured with equivalent rigour in the field.
\textcite{strohmetz2002field} randomised how restaurant servers delivered the final bill,
introducing small unexpected gifts to measure tipping lift. A single mint raised tips by
\qty{3.3}{\percent} relative to the no-gift control; two mints raised them by
\qty{14.1}{\percent}.
% strohmetz2002field: single-mint lift 3.3%, two-mint lift 14.1%, tipping field experiment
The largest effect --- approximately a \qty{23}{\percent} relative increase --- occurred
when the server delivered one mint, turned away, then returned to offer a second as a
personalised exception.
% strohmetz2002field: personalised-sequence peak ~23% relative lift
Two features of this result matter for practitioners. First, magnitude: the gift was
trivial in monetary terms; the compliance lift was not. Second, personalisation: the
perception that the concession was tailored to this specific counterpart amplified the
effect substantially. For a SaaS sales motion, a feature trial extension framed as an
exclusive, unexpected accommodation for this prospect's specific evaluation context
activates both properties simultaneously.

These heuristic cues sit at the peripheral route of the Elaboration Likelihood Model
(ELM), developed in full by \textcite{petty1986elaboration}. The ELM posits that persuasion
travels one of two routes: when a buyer has high motivation and ability to think carefully,
they follow the \emph{central route} --- weighing argument quality, scrutinising evidence,
updating beliefs rationally. When motivation or ability is low --- distracted, time-poor,
unfamiliar with the domain --- they follow the \emph{peripheral route}, relying on heuristic
cues such as descriptive norms, authority signals, or the liking they feel for the sender.
Central-route attitude change is durable and resistant to counter-argument; peripheral
attitude change fades when the cue disappears. \textcite{carpenter2015meta} synthesised 134
effect estimates from the ELM literature and confirmed this asymmetry quantitatively:
argument quality yields a substantially larger persuasion effect when elaboration is high
(central processing) than when it is low (peripheral processing).
% carpenter2015meta: 134 effects, central-route argument-quality effect >> peripheral

The practical design implication for a B2B pitch is the simultaneous satisfaction of both
routes: rigorous economic evidence for the buyer who digs into the unit economics, and
credible social-proof signals for the first thirty seconds before they decide to dig.
Generic case studies (``thousands of teams use us'') exploit social proof poorly because
the norm is too distant from the prospect's situation. The \textcite{goldstein2008field}
finding points to a concrete prescription: segment social proof to match the prospect's
exact context --- industry, company stage, regulatory environment --- and the compliance
lift follows from the provincial-norm mechanism rather than from a broader but weaker
authority halo.

Connecting influence to the acquisition funnel (\cref{eq:funnel-chain}): each stage in
the chain --- session to trial, trial to paid --- is a compliance decision. Social proof
and reciprocity are conversion-rate multipliers operating on the denominator of CAC. A
trial-to-paid conversion rate of \qty{25}{\percent} implies that three in four trials are
lost without further action; adding a contextualised social-proof signal to the trial-end
moment and a genuine reciprocity cue can shift that rate materially. The same mechanisms
apply to each funnel gate.

\begin{example}[Contextualised social proof at the trial-to-paid gate]
The SaaS funnel converts \num{10000} sessions to \num{800} trials to \num{200} paid
accounts per month, yielding a \qty{25}{\percent} trial-to-paid rate. CAC sits at
\money{600} against a target LTV:CAC of \num{5.25}\texttimes{} (\cref{eq:ltv-cac}).

The team adds two messages at the trial day-7 touchpoint: (1)~a provincial norm from
a customer in the exact segment of the prospect (``\qty{73}{\percent} of Series~B
fintech SaaS teams in your region converted to paid this quarter''), activating the
\textcite{goldstein2008field} provincial-norm mechanism; (2)~a personalised reciprocity
cue from a named account executive offering an extended onboarding session, framed as
an exception specific to this prospect's integration complexity, activating the
\textcite{strohmetz2002field} personalised-concession mechanism. Suppose a conservative
\qty{4}{\percent}point improvement in trial-to-paid, from \qty{25}{\percent} to
\qty{29}{\percent}, substantially below the field-experiment ceilings. That converts
\num{800} trials to \num{232} paid accounts --- \num{32} incremental customers --- without
changing sessions, trials, or the acquisition spend of \money{120000}. Effective CAC drops
from \money{600} to approximately \money{517}. At a CLV of \money{3150} (\cref{eq:clv}),
those \num{32} additional customers represent \money{100800} in present-value customer
equity created at zero marginal outlay. The lever is framing at a bottleneck gate, not
additional media spend.
\end{example}

\begin{admonition}{Warning}
The same mechanisms that enable ethical persuasion can be weaponised. Fake scarcity
(``only 3 left'' when the shelf is full), fabricated social proof (purchased reviews,
inflated user counts), or manufactured urgency are dark patterns: they lift the immediate
conversion metric and collapse the two things that sustain it. Trust evaporates on the
first contradiction, driving churn that destroys the CLV (\cref{eq:clv}) the conversion
was meant to create. The genre of business literature that treats manipulation as strategy
--- \textcite{greene1998laws} being one prominent example --- mistakes the short-term
tactical win for durable competitive advantage. The long-run evidence, discussed in
\cref{sec:power-trust}, runs in the opposite direction: trust compounds into LTV;
manipulation compounds into churn and reputation damage.
\end{admonition}


\section{The Psychology of the Pitch}\label{sec:pitch-psychology}
% complexity: 3

Every pitch is an exercise in framing: the same facts, presented differently, produce
different decisions. This is not a trick of rhetoric --- it is the documented structure of
human value perception. The deeper structural mechanism is the asymmetric value function
formalised in \emph{cumulative prospect theory} by \textcite{tversky1992cumulative}. The
popular treatment of this insight appears in \textcite{kahneman2011}; the primary
quantitative evidence is the 1992 paper, which derives the loss-aversion coefficient from
experimental choices over real monetary outcomes.

People do not evaluate outcomes in absolute magnitudes; they evaluate them as gains or
losses relative to a reference point, and losses weigh approximately \num{2.25} times as
heavily as equivalent gains ($\lambda \approx 2.25$, \cref{eq:prospect}).
% tversky1992cumulative: loss-aversion coefficient lambda ~= 2.25, cumulative prospect theory
The mental-accounting effects formalised by \textcite{thaler1985mental} --- and discussed
in the pricing chapter's treatment of psychological framing at \cref{sec:psych-pricing}
--- are the direct application to purchase and payment decisions: the same dollar amount
evoking loss feels worse than the same amount evoking gain, so positioning an offer as an
avoided loss rather than a captured gain increases the felt urgency to act. With
$\lambda \approx 2.25$, a pitch that frames a benefit as ``you are currently losing
\money{400}/month'' exerts more than twice the psychological gravity of a pitch that
frames the identical benefit as ``you could gain \money{400}/month''.
% tversky1992cumulative: lambda ~= 2.25 implies loss framing ~2x gain framing in persuasive gravity

The anchoring heuristic, identified by \textcite{tversky1974judgment}, is the companion
mechanism: the first number in a conversation exerts a gravitational pull on all subsequent
estimates. \textcite{galinsky2001anchoring} demonstrated in controlled experiments that
the first offer in a negotiation correlates with the final settlement price at
$r = .93$ in undefended conditions, accounting for over half --- and in some specifications
up to \qty{80}{\percent} --- of the variance in the final price.
% galinsky2001anchoring: first-offer / final-settlement r = .93, up to 80% variance explained
\textcite{northcraft1987anchoring} extended the result to domain experts: professional
real-estate agents given different listing-price anchors produced property appraisals that
tracked the anchor, even when they believed their assessments were based on objective
market analysis.
% northcraft1987anchoring: expert anchoring, real-estate professionals susceptible to listing-price anchor
In a pitch context, anchoring works before any negotiation begins: naming the cost of
inaction, or the full value of the outcome, before naming the price sets the reference
against which the price will be judged. In a negotiation context the same logic applies
within the ZOPA, treated in full in \cref{sec:negotiation}.

Dual-process theory distinguishes System~1 (fast, automatic, pattern-driven, emotionally
weighted) from System~2 (slow, deliberate, effortful, logical). System~1 fires first and
shapes the frame within which System~2 then operates. A slide deck that opens with a
visceral customer pain activates System~1 empathy before the investor's System~2 begins
evaluating the unit economics. Reversing the order --- leading with the spreadsheet ---
forces System~2 into a posture of critique before any emotional foothold exists. The
practical prescription is to sequence: hook System~1 with a vivid problem, then satisfy
System~2 with rigorous evidence. The ELM (\cref{sec:influence-principles}) explains why
both routes must be satisfied for the attitude change to persist.

For a pitch, loss-aversion translates to a disciplined reframing of the value proposition.
The question is not ``what will you gain?'' but ``what are you currently losing, and for
how long?'' The reference point is the status quo, and the status quo should cost the
counterpart something visible.

\begin{example}[Framing the enterprise tier as an avoided loss]
The SaaS offers an enterprise tier at \money{299}/month (against the standard tier at
\money{50}/month). Under a gain frame, the pitch reads: ``Upgrade to get advanced
analytics, priority support, and SSO --- value worth roughly \money{200}/month to teams
your size.'' Under a loss frame anchored to the cost of the status quo: ``Teams at your
scale without a unified analytics layer typically lose \qty{15}{\percent} of their
implementation efficiency to manual reconciliation --- at your headcount that is roughly
\money{400} of unproductive labour per month. The enterprise tier closes that gap for
\money{299}.''

The substance is identical; the reference point has shifted from zero-gain to an ongoing
loss. With $\lambda \approx 2.25$ \parencite{tversky1992cumulative}, prospect theory
(\cref{eq:prospect}) predicts that a buyer in the loss domain becomes risk-seeking ---
more willing to pay for certainty of avoiding the loss than they would be in the gain
domain for an equivalent upside. The \money{299} tier is now priced below the
\money{400} monthly pain, framing it not as a cost but as a \money{101}-per-month
recovery. System~1 closes the argument before System~2 re-opens it. To support that
System~2 close, the central route (\cref{sec:influence-principles}) demands rigorous
ROI evidence: a testimonial from a named customer with attributed efficiency outcomes
and a time-to-payback figure that a finance approver can verify.
\end{example}


\section{Negotiation: BATNA, Reservation Value \& ZOPA}\label{sec:negotiation}
% complexity: 4

Persuasion shifts beliefs; negotiation allocates the surplus that follows. The principled
negotiation framework developed at the Harvard Negotiation Project --- the popular
treatment being \textcite{fisher2011getting} --- rests on a quantitative scaffold whose
formal derivation is due to \textcite{raiffa1982art}. Its structural prescriptions ---
separate people from the problem, focus on interests not positions, invent options for
mutual gain, insist on objective criteria --- are mechanisms for reaching agreements that
are efficient (no value left on the table) and durable (both parties perceive the process
as fair, sustaining the relationship through future transactions).

Before entering any negotiation, a party must answer three questions privately. What is
my BATNA --- Best Alternative to a Negotiated Agreement, the outcome I can achieve without
this deal? What, in money terms, does my BATNA imply as a walk-away price? And where does
the counterpart's equivalent boundary lie? \textcite{kim2005power} establish from the
organisational-behaviour literature that structural power in a negotiation derives from the
strength of the BATNA: a party who can credibly walk away holds position without bluffing,
because the threat is real.
% kim2005power: BATNA strength as source of negotiation power, structural account
The monetary expression of the BATNA is the \emph{reservation value} --- the point at or
beyond which the deal is worse than no deal.

\begin{equation}\label{eq:reservation}
  RV \;=\; U(\mathrm{BATNA}),
\end{equation}
where \(U(\mathrm{BATNA})\) is the monetary utility of the best outside alternative, fully
costed (including transition costs, time, and forgone optionality). The reservation value
is not a negotiating position; it is a private calculation, fixed before the conversation
begins. Revealing it is almost always a mistake.

The \emph{Zone of Possible Agreement} (ZOPA) is the interval between the two parties'
reservation values in which both prefer a deal to their outside option:
\begin{equation}\label{eq:zopa}
  \mathrm{ZOPA} \;=\; [\,RV_{\text{seller}},\; RV_{\text{buyer}}\,],
\end{equation}
with width \(\Delta = RV_{\text{buyer}} - RV_{\text{seller}}\). If \(\Delta > 0\) a deal
is rationally possible and the negotiation is about how to split the surplus \(\Delta\).
If \(\Delta \leq 0\) no single-issue deal exists at face value, and the parties must either
introduce new variables (payment timing, volume commitments, service terms) to shift one
or both reservation values, or walk away. Every tactic --- anchoring, concession pacing,
packaging --- is ultimately a mechanism for moving the eventual settlement point within
the ZOPA toward the tactician's reservation value.

The anchoring results of \textcite{galinsky2001anchoring} bear directly on the question of
who makes the first offer inside the ZOPA. The high $r = .93$ correlation between first
offer and final settlement in undefended conditions means that the party who anchors first
--- at an ambitious but ZOPA-respecting value --- structurally advantages itself.
% galinsky2001anchoring: first-offer anchor advantage applies within the ZOPA
\textcite{northcraft1987anchoring} confirm that even expert negotiators absorb the anchor
and adjust insufficiently; the defence, also identified by \textcite{galinsky2001anchoring},
is to focus exclusively on one's own target price and the counterpart's BATNA rather than
engaging with the anchor as a reference point.
% galinsky2001anchoring: focusing on own target / counterpart BATNA nullifies anchor effect

\begin{example}[Enterprise deal and a Series-A term sheet]
\textbf{Enterprise deal.} The SaaS is negotiating an enterprise contract. The seller's
BATNA is to close three additional mid-market accounts at \money{18000} each (total
\money{54000} annually), with lower support overhead; \(RV_{\text{seller}} = \money{60000}\)
after netting transition costs. The enterprise buyer's BATNA is a custom internal build
estimated at \money{110000} in engineering and maintenance;
\(RV_{\text{buyer}} = \money{100000}\) after a \qty{10}{\percent} risk haircut on build
uncertainty. The ZOPA is \([\money{60000}, \money{100000}]\), width
\(\Delta = \money{40000}\). The seller opens at \money{95000} --- an anchor near the
buyer's ceiling, grounded in stated platform value and named comparable deployments ---
and the negotiation settles at \money{82000}: within the ZOPA, above both reservation
values, and sustainable for both. The \money{22000} surplus above the seller's reservation
value and the \money{18000} surplus below the buyer's reflect the anchoring advantage of
the first move and the concession discipline of both parties.

\textbf{Series-A term sheet.} The same logic applies to equity financing
(\cref{ch:fundraising-and-dilution}). Suppose the founder's BATNA is a competing term
sheet at a \money{12000000} pre-money valuation. \(RV_{\text{founder}}\) is the dilution
implied by that valuation plus the cost of delaying. The lead investor's BATNA is
deploying capital in a comparable company at a \money{15000000} pre-money; their
\(RV_{\text{investor}}\) is the maximum pre-money at which the deal still meets their
return hurdle. The ZOPA is the valuation band between those two thresholds. The founder
who anchors first with a specific pre-money figure supported by a named comparable ---
objective-criteria-backed, per \textcite{raiffa1982art} --- pulls the settlement toward
their reservation value. The investor who delays revealing their portfolio-wide hurdle
rate protects theirs. Both are operating within the principled-negotiation framework.
\end{example}

\begin{admonition}{Warning}
Two failure modes are common. First, anchoring beyond the counterpart's reservation
value --- opening so aggressively that the other party exits rather than counters ---
collapses the ZOPA by signalling bad faith. The anchor must be ambitious but plausible,
grounded in objective criteria. Second, the winner's curse: in competitive bidding, the
winner is systematically the party that most overestimated the asset's value
\parencite{thaler1988winnerscurse}. A founder who ``wins'' a hire by bidding far above
the candidate's reservation value has overpaid relative to market; a venture firm that
wins a deal by bidding far above its return hurdle has destroyed expected return.
% thaler1988winnerscurse: winner's curse — winning bidder overestimates asset value
Both errors are cured by anchoring the reservation value to the BATNA's monetary utility,
not to desire.
\end{admonition}


\section{Power \& Sustainable Influence}\label{sec:power-trust}
% complexity: 3

The negotiation frameworks of \cref{sec:negotiation} assume that the parties meet at the
table. Before the table, power determines whether the meeting happens, what agenda gets
set, and whose framing shapes the opening anchor. \textcite{pfeffer2010power} provides the
evidence-based organisational account: power accrues to those who control critical resource
dependencies, build broad coalitions, and demonstrate competence visibly. The contrast with
the popular mythology of the genre is worth naming briefly. \textcite{greene1998laws}
constructs a taxonomy of Machiavellian tactics from historical anecdotes; it is not
derived from systematic evidence. The strategies it endorses --- concealment, strategic
dependence-creation --- are precisely those that destroy the trust required to retain
employees, partners, and customers over time. Pfeffer's structural account describes
dynamics that compound; anecdotal manipulation tactics erode.

The durable foundation of sustainable influence is trust. \textcite{maister2000trusted}
decompose trust into its structural components as a practitioner heuristic:
\begin{equation}\label{eq:trust}
  T \;=\; \frac{C + R + I}{S},
\end{equation}
where \(C\) is credibility (the quality of what you say), \(R\) is reliability (the
consistency between what you say and what you do), \(I\) is intimacy (the degree to which
the counterpart believes you will handle sensitive information with discretion), and \(S\)
is self-orientation (the degree to which the counterpart perceives your interest as being
in yourself rather than in them). The structure of the equation is the operational insight:
\(S\) is in the denominator, meaning that even high credibility and reliability collapse
if the counterpart reads the advisor as self-serving.

The practitioner heuristic of \cref{eq:trust} is grounded in the integrative organisational
model of \textcite{mayer1995model}, who formalise trustworthiness as an additive structure
of three independently observable dimensions: ability (domain-specific competence),
benevolence (the perception that the trustee desires the trustor's success, not merely their
own), and integrity (word--deed consistency and perceived adherence to acceptable principles).
\textcite{colquitt2007meta} provide the meta-analytic validation across 132 independent
samples: the corrected correlation between each dimension and the resulting level of trust
is \(r_c = .67\) for ability, \(r_c = .63\) for benevolence, and \(r_c = .62\) for
integrity.
% colquitt2007meta: ability r_c=.67, benevolence r_c=.63, integrity r_c=.62, N~6800–7300 per dimension
The three components contribute with near-equal weight --- no single dimension dominates ---
which matches the balanced numerator structure of \cref{eq:trust}. Critically, trust
translates directly into the risk-taking behaviour that makes a procurement officer sign a
multi-year contract: the corrected correlation between trust and risk-taking is
\(r_c = .42\) (\(N = 1{,}384\)) \parencite{colquitt2007meta}.
% colquitt2007meta: trust -> risk-taking r_c=.42, N=1384

Reducing self-orientation --- asking questions rather than pitching, disclosing conflicts
proactively, occasionally directing a client to a competitor who is a better fit ---
multiplies the trust numerator without adding to it. This is not mere altruism: a customer
who enters a contract with high trust (\(T\)) exhibits lower churn (the monthly churn rate
of \qty{0.65}{\percent} assumed in the running model is predicated on customers remaining
convinced of the vendor's benevolence through each renewal cycle), higher willingness to
expand (moving from the standard \money{50}/month tier toward the enterprise tier), and
lower CAC through referrals that originate inside the trust network rather than from paid
acquisition (\cref{eq:cac}). The ROIC implication (\cref{eq:roic}) follows directly: trust
lowers the CAC numerator, lowers churn-driven attrition in the CLV denominator, and widens
margin --- each of which raises the return on invested capital.

\begin{admonition}{Contested}
\Cref{eq:trust} is a practitioner heuristic, not an estimated regression equation.
\textcite{maister2000trusted} derive it from consulting practice, not from a controlled
empirical study. Its virtue is structural --- it correctly identifies \(S\) as the
multiplicative divisor --- but the precise ratio is not empirically calibrated. The
underlying science is \textcite{mayer1995model} and \textcite{colquitt2007meta}; the
equation is the mnemonic that makes the science operational. Use the component correlations
from \textcite{colquitt2007meta} to prioritise where to invest: because ability, benevolence,
and integrity contribute nearly equally, a deficit in any single dimension is not
compensated by surplus in the others.
\end{admonition}

\begin{example}[Founder trust as a compounding asset]
Consider a SaaS founder building investor and talent relationships over two years before
a Series-A raise (\cref{ch:fundraising-and-dilution}). The founder takes four visible
steps, each corresponding to a component of \cref{eq:trust}:

\begin{itemize}
  \item \textbf{Credibility (\(C\)):} publishes monthly metrics updates to a small list
    of angel investors and operators, including months when growth was below target, with
    honest attribution of root cause. This builds the \emph{ability} signal of
    \textcite{mayer1995model} through demonstrated domain competence and the integrity
    signal through word--deed consistency.
  \item \textbf{Reliability (\(R\)):} delivers every commitment made to early investors
    --- product roadmap dates, hiring milestones, callback timelines --- without exception
    over eight quarters. The \(r_c = .62\) integrity channel \parencite{colquitt2007meta}
    compounds with each kept promise.
    % colquitt2007meta: integrity r_c=.62, trust antecedent
  \item \textbf{Intimacy (\(I\)):} advises two prospective hires, in confidence, not to
    join the company yet because the role is not well-defined; both join one year later
    when the role is real. This demonstrates benevolence (the \(r_c = .63\) channel)
    at a visible cost to short-term hiring velocity.
    % colquitt2007meta: benevolence r_c=.63, trust antecedent
  \item \textbf{Self-orientation (\(S\)):} in investor meetings, opens with questions
    about the firm's current portfolio concerns before describing the company's needs.
\end{itemize}

The measurable consequences appear at the raise. Referral-sourced inbound from the trust
network costs zero in paid CAC (\cref{eq:cac}). Two of the three Series-A investors the
founder most wanted take meetings without a warm introduction, because the credibility
record is already legible. The lead closes at a pre-money valuation in the upper half of
the modelled range, partly because investors with high \(T\) scores command a lower
uncertainty discount. Three of the six key engineering hires come through talent referrals
from the candidate the founder previously advised away --- trust converting into recruitment
CAC elimination. The compounding mechanism is the same as in \cref{eq:clv}: each period
of trust-reinforcing behaviour raises the relationship's net present value by lowering
the implicit discount rate the counterpart applies to future interactions.
\end{example}
